\documentclass[10pt]{article}
\usepackage[margin=1in]{geometry}
\usepackage{amsmath,amssymb}
\usepackage{booktabs}
\usepackage{graphicx}
\usepackage{xcolor}
\usepackage[hidelinks]{hyperref}
\title{Failure-Directed Domain Randomization for Cooperative Robot-Soccer Skill Selection}
\author{Research software report}
\date{\today}

\begin{document}
\maketitle

\begin{abstract}
We present a compact multi-fidelity testbed for parameter-shared multi-agent reinforcement
learning in a two-attacker robot-soccer task. Policies train in a fast abstract simulator and are
frozen before zero-shot evaluation in an independently implemented Pymunk rigid-body simulator.
The proposed curriculum directs domain randomization toward perturbation profiles on which the
current policy fails. Pymunk transfer is sim-to-sim evaluation and a controlled proxy for selected
sim-to-real robustness questions, not physical deployment. An initial Phase 1 audit found a
saturated transfer benchmark and metric defects. The repaired implementation now passes explicit
boundary, aggregation, paired-seed, and policy-sensitivity safeguards. A clean coherent rerun now
passes the Phase 2 launch gate while exposing a 0.530 nominal cross-simulator success gap.
\end{abstract}

\section{Introduction and motivation}
Robot soccer combines partial observability, delayed control, contact dynamics, communication
failure, and cooperation. End-to-end physical training is expensive, while conclusions drawn from
one simulator can be brittle. We therefore ask whether a decentralized, cooperative skill policy
learned in a deliberately simple simulator survives a change in simulator fidelity and unseen
perturbations. The contribution is a failure-directed curriculum that spends training episodes on
the policy's empirically weakest named perturbation profiles.

\section{RoboCup and robot-soccer context}
RoboCup established robot soccer as a long-horizon benchmark joining perception, control, and
multi-agent reasoning \cite{kitano1997robocup}. Modern learned soccer controllers demonstrate the
value of reinforcement learning for dynamic bipedal play \cite{haarnoja2023soccer}. Our scope is
narrower: high-level behavior selection for two abstract attackers, assuming lower-level motion
and perception stacks.

\section{Related work}
PPO uses a clipped surrogate objective to stabilize on-policy updates
\cite{schulman2017ppo}. MAPPO showed that careful centralized-value training can be effective in
cooperative multi-agent tasks \cite{yu2022mappo}; this repository uses centralized training with
decentralized execution. Domain randomization varies simulator parameters to improve transfer
\cite{tobin2017domain}, while automatic curricula adapt the distribution as learning proceeds
\cite{openai2019rubik}. The present method uses explicit, interpretable failure profiles rather
than an unstructured Cartesian parameter distribution. Cross-fidelity evaluation follows the
broader motivation of using simulation variation to reveal model-specific policy behavior.

\section{Problem formulation}
We model the task as a cooperative partially observable stochastic game
$\langle\mathcal{S},\{\mathcal{O}_i\},\{\mathcal{A}_i\},P,r,\gamma\rangle$.
Two agents receive distinct egocentric observations $o_i$, select discrete skills through one
shared policy $\pi_\theta(a_i\mid o_i)$, and receive one team reward. During MAPPO training only, a
critic $V_\phi(s)$ observes the true joint state. The actor remains decentralized at evaluation.

\section{Abstract simulator}
The training simulator explicitly integrates bounded acceleration, speed, turn rate, ball drag,
collisions, boundary and goal events, kick impulses, a delayed-action queue, perception queues, and
imperfect teammate messages. Each policy step repeats a selected macro-action for four low-level
steps. The scripted defender can guard, chase, or predict an interception point and clears the ball
through ordinary limited motion rather than teleportation.

\section{Pymunk transfer simulator}
The transfer environment rebuilds players and the ball as Pymunk rigid bodies. Velocity targets
are realized by bounded forces and torque; kicks set a target ball velocity through a mass-scaled
impulse delivered at contact range. The ball crosses field lines and is classified by the same
goal, clear, and out-of-bounds rules as the abstract simulator, while player-only static boundaries
and numerical clipping keep every player in the playable rectangle. Inertia, collision response,
ball drag, integration substeps, and reaction delay still differ from the abstract equations.
Observation, action, reward, and metric interfaces remain fixed, so the actor transfers unchanged.
This is zero-shot cross-simulator transfer, not completed physical sim-to-real deployment.

\section{Observation and action design}
Each 61-dimensional local vector contains self motion and action state, delayed ball history,
egocentric goalposts and boundaries, teammate kinematics and stale message state, defender state,
confidence flags, and identity. The seven actions are approach, left/right dribble, shoot, pass,
move to support, and hold/face ball. Noise affects perceived state only; true dynamics and
centralized state remain separate.

\section{IPPO and MAPPO}
Both algorithms use one categorical actor. IPPO shares a local value network and computes one
advantage per local trajectory. MAPPO uses one centralized team value per environment step and
broadcasts the team advantage to the two actor samples; its value loss is counted once. GAE
bootstraps through time-limit truncation but does not propagate across a reset. Separate Adam
optimizers, clipped policy and value objectives, entropy regularization, gradient clipping,
target-KL stopping, and linear learning-rate decay follow standard PPO practice.

\section{Failure-directed domain randomization}
Let $f_i$ be a smoothed failure rate for profile $i$. After warm-up, deterministic validation on
disjoint seeds updates
\begin{align}
 w_i &= (\epsilon + f_i)^\alpha,\\
 \tilde p_i &= (1-\beta)\frac{w_i}{\sum_j w_j}+\frac{\beta}{N}.
\end{align}
The distribution is projected onto the simplex with a hard maximum profile probability. The
uniform mixture preserves exploration; the cap prevents one persistent failure from consuming all
training. Curriculum evaluation never updates observation normalization.

\section{Experimental protocol}
Training, ordinary validation, curriculum validation, abstract test, Pymunk transfer, and video
seeds occupy disjoint ranges. Policies are selected on fixed abstract validation seeds, frozen,
then evaluated on nominal defenders, all profiles, and an action-delay by localization-noise grid.
Multiple training seeds are aggregated by experiment name. Checkpoints include all optimizer,
normalizer, RNG, schedule, and curriculum state needed to resume.

\section{Baselines and ablations}
The matrix includes random, double-chase, and hysteretic role-based heuristics; nominal IPPO;
nominal MAPPO; uniformly randomized MAPPO; and failure-directed MAPPO. Focused ablations remove
action-delay, communication, or observation-noise profile families while retaining the same code
path and configuration inheritance.

\section{Metrics}
Primary measures are goal success, return, time to score, and the abstract-to-Pymunk success gap.
Secondary measures include pass completion, possession, redundant chasing, collisions, invalid
skills, action switches, worst-decile return, CVaR, explicitly separated defender-mode and
perturbation-profile minima, normalized robustness area, and bootstrap intervals. Reward-component
logs distinguish sparse task reward from small
potential-based progress, possession, and support shaping.

\section{Results}
\IfFileExists{reports/generated_results.tex}{% Phase 1 status after local repair validation plus frozen-checkpoint re-evaluation.
% scripts.compare_runs may replace this file after the Phase 2 multi-seed study.
\newcommand{\ResultsAvailable}{phaseone}
\newcommand{\ResultsStatus}{The repaired Phase 1 benchmark passes its local readiness audit. One clean post-repair notebook run is required for provenance before Phase 2; comparative MAPPO results remain pending.}
}{
  \newcommand{\ResultsStatus}{Final empirical results have not yet been generated.}}
\begin{center}
\fcolorbox{black}{gray!10}{\parbox{0.88\linewidth}{\centering\ResultsStatus}}
\end{center}

\subsection{Phase 1 clean rerun and readiness audit}
The clean post-repair notebook run was produced on 21 July 2026 at commit \texttt{16c94cd}. The
fresh IPPO actor completed 2,007,040 abstract-simulator steps with no recorded failure. Mean
training success increased from 0.088 over the first 25 updates to 0.837 over the last 25. The
selected checkpoint was saved at 1,277,952 steps, where a fixed 40-episode validation set reached
0.975 success; final validation was 0.950. These are credible signs that the implementation can
optimize the nominal training task, but they remain one-seed diagnostics rather than comparative
robustness evidence.

\begin{table}[ht]
\centering
\small
\begin{tabular}{lllrrr}
\toprule
Method & Simulator & Defender scope & Episodes & Success & Mean return \\
\midrule
Random & abstract & intercept & 100 & 0.030 & $-4.585$ \\
Double chase & abstract & intercept & 100 & 0.150 & $-3.624$ \\
Role based & abstract & intercept & 100 & 0.280 & 0.682 \\
Nominal IPPO & abstract & all three & 300 & 0.417 & 1.653 \\
Nominal IPPO & abstract & intercept & 100 & 0.870 & 7.722 \\
Nominal IPPO & abstract & stationary goalie & 100 & 0.330 & 0.392 \\
Nominal IPPO & abstract & chase ball & 100 & 0.050 & $-3.153$ \\
Random & Pymunk & intercept & 100 & 0.490 & 0.242 \\
Double chase & Pymunk & intercept & 100 & 0.680 & 5.526 \\
Role based & Pymunk & intercept & 100 & 0.660 & 5.884 \\
Frozen nominal IPPO & Pymunk & intercept & 100 & 0.340 & 0.247 \\
\bottomrule
\end{tabular}
\caption{Post-repair Phase 1 diagnostics on new fixed test seeds: 130,000--130,099 in abstract and
140,000--140,099 in Pymunk. All rows come from the clean post-repair notebook; the Pymunk row
evaluates the same frozen best actor without updating its weights. ``All three'' averages equal-sized
stationary-goalie, chase-ball, and intercept subsets and should not be compared directly with a
single-defender baseline.}
\label{tab:phase1-checkpoint}
\end{table}

The nominal policy solves its trained intercept opponent substantially more often than every
scripted abstract baseline. It does not demonstrate broad opponent robustness: success is 0.050
against the chase-ball defender and 0.330 against the stationary goalie. Across the 300 abstract
and 100 Pymunk learned-policy diagnostic episodes it makes no valid pass attempt, so Phase 1 still
does not show that success depends on learned passing cooperation.

The repaired Pymunk benchmark is no longer saturated. Scripted success spans 0.490--0.680, and the
role policy reduces redundant chasing from 0.475 for double-chase to 0.029 while retaining 0.660
success. On paired seeds 140,000--140,099, fresh IPPO succeeds on 34 episodes versus 49 for random:
18 seeds are successes for both, 35 failures for both, 16 IPPO-only successes, and 31 random-only
successes. A two-sided exact paired test on the 47 discordant seeds gives $p=0.040$. Thus the clean
nominal actor transfers significantly worse than random under this paired test and remains below
both deliberate heuristics. This is a policy-transfer failure rather than a saturated-benchmark
artifact. Matched intercept success falls from 0.870 in abstract dynamics to 0.340 in Pymunk, a
0.530 cross-simulator gap that directly motivates Phase 2's robustness methods.

The instrumentation defects are corrected. Possession receives one sample per macro step;
possession, redundant-chase, and invalid-action fractions are range checked; invalid attempts use
the two-agent action count as their denominator; defender-mode minima are explicit; and profile
minima are emitted only for multi-profile suites. The observed defender-mode minimum is 0.050,
which honestly exposes the frozen actor's chase-ball weakness.

\subsection{Phase 2 launch decision}
The clean post-repair run produces a \textbf{go}. Every coded
gate passes: six baseline cells each contain 100 episodes, all reported fractions are bounded,
random Pymunk success is below 0.50, the Pymunk method spread exceeds 0.15, role coordination beats
random and reduces redundant chasing, all three defender modes have 100 learned-policy episodes,
and nominal IPPO clearly exceeds the scripted abstract controls. The baseline and fresh nominal
actor, evaluations, videos, resolved configurations, commit, and
\texttt{phase1\_readiness.json} now belong to one coherent post-repair run, and that file reports
\texttt{phase2\_ready: true}. This closes the original no-go. The Phase 2 notebook independently
recomputes the same checks and refuses to launch MAPPO otherwise.

This is an engineering launch gate, not a robustness claim. Random Pymunk success is 0.490, only
0.010 below the saturation threshold, and its 95\% bootstrap interval extends above 0.50. More
importantly, the nominal actor's paired Pymunk result is worse than random. Phase 2 should proceed
because these weaknesses are now measurable on a policy-sensitive benchmark, not because Phase 1
has solved transfer.

The comparison command should replace generated content only from completed, validated artifacts.
Smoke-run observations remain infrastructure diagnostics and must not be interpreted as final
evidence.

\IfFileExists{reports/figures/method_success_comparison.png}{
\begin{figure}[t]\centering
\includegraphics[width=0.78\linewidth]{reports/figures/method_success_comparison.png}
\caption{Observed method comparison generated from registered artifacts.}
\end{figure}}{}

\section{Qualitative behavior and failure analysis}
The video workflow records nominal, transfer, and forced-profile trajectories with action and
outcome overlays. Analysis should distinguish coordination failures (role oscillation, redundant
chasing, poor support), control failures (delay and movement response), perception/communication
failures, and contact-model failures. Representative examples should use fixed seeds rather than
favorable outcomes. The clean run records three nominal IPPO videos in each simulator, including
successes and out-of-bounds failures. A separate post-repair spot check recorded Pymunk role-policy
successes and an out-of-bounds failure with outcome overlays; players remained inside the rendered
field and labels matched terminal events. These are mechanics checks and qualitative examples,
not representative behavioral rates.

\section{Limitations and claim boundaries}
The task omits images, recurrent memory, joint-level locomotion, self-play, learned opponents, and
full-team tactics. The scripted skill adapters still encode useful soccer structure. Pymunk does
not model sensing, compliance, actuator limits, falls, or hardware damage. Consequently the study
is multi-fidelity and sim-to-real-oriented; it cannot establish physical transfer or safety.

\section{Conclusion}
The repository supplies an auditable experiment for asking whether adaptive, failure-directed
variation improves cooperative policy robustness across simulator implementations. The repaired
Phase 1 benchmark is policy-sensitive, its metrics are bounded and correctly grouped, and the
frozen nominal actor still learns the trained abstract task while exhibiting a measurable Pymunk
gap. Phase 2 is technically ready now that the clean Phase 1 provenance run is complete; comparative robustness
claims remain pending the planned multi-seed study. A later Webots or robot phase should be
attempted only after the simulator protocol is reproduced and a complete perception, localization,
locomotion, and safety stack is available.

\bibliographystyle{plain}
\bibliography{reports/references}
\end{document}
